\section{KHAI PHÁ DỮ LIỆU}

\subsection{Khái niệm khai phá dữ liệu}
Khai phá dữ liệu là quá trình khai phá và phát hiện các mối tương quan, mẫu và xu hướng mới có ý nghĩa bằng cách phân tích một lượng lớn dữ liệu được lưu trữ, sử dụng các kỹ thuật nhận dạng mẫu và các phương pháp thống kê, toán học.

\subsection{Các lĩnh vực liên quan}
\begin{enumerate}
    \item \textbf{Thống kê :} Nền tảng toán học của khai phá dữ liệu.
    \begin{itemize}
        \item Vai trò: Cung cấp các công cụ để mô tả dữ liệu, phân tích sự phụ thuộc và đưa ra kết luận từ dữ liệu.
        \item Điểm khác biệt: Thống kê thường tập trung vào việc kiểm định giả thuyết trên các tập dữ liệu nhỏ hơn và sạch hơn. Khai phá dữ liệu sử dụng thống kê để tìm ra các mẫu trong dữ liệu hỗn loạn và khổng lồ.
        \item Các khái niệm chung: Hồi quy, Phân phối chuẩn, Phương sai, Độ lệch chuẩn.
    \end{itemize}

    \item \textbf{Trí tuệ nhân tạo \& Học máy :} Cung cấp các thuật toán để khai phá.
    \begin{itemize}
        \item Vai trò: Cung cấp các kỹ thuật học và tự cải thiện từ dữ liệu mà không cần lập trình cụ thể cho từng trường hợp.
        \item Điểm khác biệt:
        \begin{itemize}
            \item \textit{Machine Learning:} Tập trung vào việc tối ưu hóa độ chính xác của dự đoán.
            \item \textit{Data Mining:} Tập trung vào việc tìm ra các tính chất chưa biết hoặc các mẫu ẩn.
        \end{itemize}
        \item Các kỹ thuật chung: Mạng nơ-ron, Decision Trees, thuật toán di truyền.
    \end{itemize}

    \item \textbf{Hệ cơ sở dữ liệu :} Cung cấp môi trường lưu trữ và quản lý dữ liệu.
    \begin{itemize}
        \item Vai trò: Lưu trữ, truy xuất và quản lý dữ liệu hiệu quả cho khai phá dữ liệu.
        \item Các khái niệm liên quan:
        \begin{itemize}
            \item \textit{OLTP (Online Transaction Processing):} Xử lý giao dịch hàng ngày.
            \item \textit{OLAP (Online Analytical Processing):} Xử lý phân tích trực tuyến.
            \item \textit{Data Warehousing:} Kho dữ liệu tập trung đã được làm sạch để phục vụ phân tích.
        \end{itemize}
    \end{itemize}

    \item \textbf{Trực quan hóa dữ liệu :} Cầu nối giữa kết quả phân tích và người sử dụng.
    \begin{itemize}
        \item Vai trò: Thể hiện mẫu, xu hướng và quan hệ bằng biểu đồ.
        \item Công cụ: Tableau, Power BI, Matplotlib.
    \end{itemize}

    \item \textbf{Các lĩnh vực liên quan khác:}
    \begin{itemize}
        \item Nhận dạng mẫu : Phân loại dữ liệu đầu vào dựa trên các đặc trưng đã biết.
        \item Truy hồi thông tin : Tìm kiếm tài liệu và thông tin trong các tập dữ liệu phi cấu trúc.
        \item Tính toán hiệu năng cao : Cung cấp xử lý song song, phân tán cho quy mô lớn.
    \end{itemize}
\end{enumerate}

\subsection{Các nhiệm vụ của Khai phá dữ liệu}
Các nhiệm vụ chính của khai phá dữ liệu bao gồm:
\begin{itemize}
    \item Phân lớp
    \item Phân cụm
    \item Phát hiện bất thường
    \item Khai phá luật kết hợp
    \item Hồi quy/dự báo
\end{itemize}

\subsection{Ứng dụng thực tế}
Khai phá dữ liệu được áp dụng rộng rãi trong nhiều lĩnh vực đời sống:
\begin{itemize}
    \item Tài chính - Ngân hàng: Phát hiện gian lận và giao dịch lừa đảo; Dự báo rủi ro tín dụng.
    \item Thương mại điện tử - Marketing: Hệ thống gợi ý sản phẩm, phân khúc khách hàng.
    \item Y tế: Dự đoán bệnh lý; Hỗ trợ chẩn đoán qua phân tích ảnh y khoa.
\end{itemize}

\subsection{Vai trò trong Công nghiệp 4.0}
Công nghiệp 4.0 được xây dựng trên nền tảng của Big Data + AI + IoT. Trong đó, Khai phá dữ liệu đóng vai trò là xương sống với các nhiệm vụ:
\begin{itemize}
    \item Khai thác dữ liệu khổng lồ từ các cảm biến IoT.
    \item Hỗ trợ ra quyết định tự động hóa.
    \item Tối ưu hóa quy trình sản xuất.
    \item Phát hiện sự cố và rủi ro trong thời gian thực.
\end{itemize}

\subsection{Phân biệt Khai phá dữ liệu và Tìm kiếm thông thường}
\begin{table}[H]
    \centering
    \renewcommand{\arraystretch}{1.3}
    \begin{tabularx}{\textwidth}{|p{3cm}|X|X|}
        \hline
        \textbf{Tiêu chí} & \textbf{Tìm kiếm thông thường} & \textbf{Khai phá dữ liệu} \\
        \hline
        \textbf{Mục tiêu} & Tìm thông tin cụ thể, chính xác đã tồn tại. & Tìm các mẫu (patterns), tri thức mới tiềm ẩn. \\
        \hline
        \textbf{Kết quả} & Dữ liệu, tài liệu khớp với từ khóa truy vấn. & Mô hình, quy luật, xu hướng, mối liên hệ. \\
        \hline
    \end{tabularx}
\end{table}